%=======================================================================
%  Care-Sync / M-Health : Full Technical Report
%  Personalized Physiological Forecasting & Anomaly Detection Platform
%
%  Compile with:  pdflatex Care-Sync_Technical_Report.tex   (run twice for ToC)
%=======================================================================
\documentclass[11pt,a4paper]{article}

\usepackage[utf8]{inputenc}
\usepackage[T1]{fontenc}
\usepackage[margin=1in]{geometry}
\usepackage{amsmath,amssymb,amsfonts}
\usepackage{booktabs}
\usepackage{longtable}
\usepackage{array}
\usepackage{enumitem}
\usepackage{xcolor}
\usepackage{listings}
\usepackage{graphicx}
\usepackage{caption}
\usepackage{titlesec}
\usepackage[hidelinks]{hyperref}

\definecolor{codebg}{RGB}{245,247,249}
\definecolor{codekw}{RGB}{0,102,153}
\definecolor{codecm}{RGB}{110,120,130}
\definecolor{codest}{RGB}{170,55,55}
\definecolor{emerald}{RGB}{4,120,87}

\lstset{
  basicstyle=\ttfamily\footnotesize,
  backgroundcolor=\color{codebg},
  keywordstyle=\color{codekw}\bfseries,
  commentstyle=\color{codecm}\itshape,
  stringstyle=\color{codest},
  breaklines=true,
  showstringspaces=false,
  frame=single,
  rulecolor=\color{gray!40},
  columns=fullflexible,
  language=Python,
  numbers=none,
  aboveskip=8pt,belowskip=8pt
}

% inline code helper that tolerates underscores
\newcommand{\code}[1]{\lstinline[basicstyle=\ttfamily\normalsize]!#1!}

\titleformat{\section}{\Large\bfseries\color{emerald}}{\thesection}{0.7em}{}
\titleformat{\subsection}{\large\bfseries}{\thesubsection}{0.6em}{}

\title{\textbf{Care-Sync / M-Health}\\[4pt]
\large A Privacy-First Platform for Personalized Physiological Forecasting\\
and Explainable Anomaly Detection\\[6pt]
\normalsize Full Technical Report --- Data Pipeline, Feature Engineering,\\
Models, Anomaly Layer, and System Architecture}
\author{Care-Sync Engineering}
\date{\today}

\begin{document}
\maketitle
\begin{abstract}
\noindent
Care-Sync is an end-to-end platform that ingests wearable physiological signals
(Empatica~E4: PPG/BVP, HR, inter-beat intervals, accelerometry, electrodermal
activity, and skin temperature), learns what is \emph{normal for a specific
individual}, forecasts their heart rate and heart-rate variability (HRV) one,
five, and ten minutes ahead, and raises \emph{explainable} health-drift alerts
when the measured signal deviates from that personal baseline. This report
documents every step of the working system: how raw device files are read,
synchronized, filtered, and cleaned; how each feature is extracted (with exact
formulas and windowing); how sequence windows and forecast targets are
constructed; the forecasting models (TCN, LSTM, GRU, Transformer, XGBoost) and
their heteroscedastic $\beta$-NLL training; the personalization stack (global
pretraining, per-user fine-tuning, self-supervised warm-start, online updates,
and the digital twin); the multi-channel anomaly detector (adaptive z-scores,
Kalman innovation, CUSUM drift, illness composite, signal-quality gating, and
hysteresis alerting); and the surrounding platform (FastAPI AI service, Django
REST backend, and Next.js frontend).
\end{abstract}

\tableofcontents
\newpage

%=======================================================================
\section{System Overview}
%=======================================================================

Care-Sync is organized into four cooperating layers plus a device-side helper,
each in its own top-level directory of the repository.

\begin{longtable}{@{}p{0.20\textwidth}p{0.14\textwidth}p{0.58\textwidth}@{}}
\toprule
\textbf{Layer} & \textbf{Stack} & \textbf{Responsibility} \\
\midrule
\endhead
\code{frontend/} & Next.js 16, React 19, TypeScript & Web UI: login, dashboard,
files, PPG visualization, the device \emph{Signal Portal}, trends, alerts,
consent, admin. Talks to the backend over \code{/api}. \\
\code{backend/} & Django 6, DRF, SimpleJWT & REST API: authentication (email
OTP + JWT), file management, signal visualization, HL7 export, device registry,
portal aggregation, alert rules, HRV anomaly records, consent. \\
\code{ai/} & FastAPI, LangGraph, PyTorch & Agentic analysis + the HRV
forecast/anomaly/digital-twin service. Wraps the research pipeline behind a
stable HTTP contract with a deterministic mock fallback. \\
\code{HRV\ /personalized\_hrv\_system/} & Python, NumPy, pandas, SciPy, PyTorch,
XGBoost & The research pipeline: cleaning, feature engineering, models,
training, personalization, anomaly scoring, simulation, evaluation. \\
\code{device-ingestion/} & stdlib Python & Device heartbeat helper that posts
battery/firmware/session status to the backend device API. \\
\bottomrule
\end{longtable}

\subsection{End-to-end data flow}

The physiological data path, from sensor to alert, is:

\begin{enumerate}[leftmargin=1.4em,itemsep=2pt]
\item \textbf{Ingest} raw per-signal files, each carrying its own start time and
      sample rate in a header (Section~\ref{sec:ingest}).
\item \textbf{Synchronize} every signal onto a common 1\,Hz grid, aggregating
      fast signals per second (Section~\ref{sec:sync}).
\item \textbf{Clean} each modality with signal-appropriate filters, artifact
      rejection, and gap handling (Section~\ref{sec:clean}).
\item \textbf{Engineer features}: rolling HR/HRV statistics, activity, circadian
      baseline, physiological state, recovery, subject-relative baselines, and
      long-context summaries (Section~\ref{sec:features}).
\item \textbf{Window} the feature table into input sequences with multi-horizon
      forecast targets (Section~\ref{sec:windows}).
\item \textbf{Forecast} HR (and RMSSD) at $+60$, $+300$, $+600$\,s with a
      predicted mean \emph{and} uncertainty (Section~\ref{sec:models}).
\item \textbf{Personalize}: bootstrap from a population model, fine-tune per
      user, and keep adapting online; summarize the person as a digital twin
      (Section~\ref{sec:personal}).
\item \textbf{Detect anomalies} by comparing actual vs.\ forecast through a
      multi-channel, quality-gated, context-aware scorer, and raise explainable
      alerts (Section~\ref{sec:anomaly}).
\item \textbf{Serve} all of this over HTTP and surface it in the portal
      (Sections~\ref{sec:aiservice}--\ref{sec:frontend}).
\end{enumerate}

A central design principle is \textbf{causality}: every feature and every
baseline is computed using only past and current samples (rolling windows,
expanding means with a one-step shift, EWMA trackers), so the exact same code
that trains offline can run in real time on a streaming device.

%=======================================================================
\section{Data Ingestion}\label{sec:ingest}
%=======================================================================

Raw data follows the Empatica~E4 export format. Each session is a folder
(\code{backend/Users/<user>/<session>/} in production, or
\code{Dataset3/Raw\_data/SXX/} in research) containing one file per signal.
Readers live in \code{src/data/loader.py}.

\begin{longtable}{@{}llp{0.52\textwidth}@{}}
\toprule
\textbf{File} & \textbf{Native rate} & \textbf{Format} \\
\midrule
\endhead
\code{HR.csv} & $\sim$1\,Hz & row1 = Unix start ts (UTC), row2 = Hz, then values. \\
\code{EDA.csv} & 4\,Hz & same 2-row header, electrodermal activity ($\mu$S). \\
\code{TEMP.csv} & 4\,Hz & same header, skin temperature ($^\circ$C). \\
\code{BVP.csv} & 64\,Hz & same header, raw blood-volume-pulse (PPG) waveform. \\
\code{ACC.csv} & 32\,Hz & row1 = start ts $\times3$, row2 = Hz $\times3$, then
$x,y,z$ in units of $1/64\,g$ (converted to $g$ on read). \\
\code{IBI.csv} & irregular & row0 = \texttt{start\_ts, IBI}; then
$(\text{offset}_s, \text{ibi}_s)$ pairs --- inter-beat intervals, no fixed Hz. \\
\code{tags*.csv} & --- & one Unix timestamp per line (event/button markers). \\
\code{info.txt} & --- & metadata, shown verbatim. \\
\bottomrule
\end{longtable}

Each reader parses the header to build a UTC \code{DatetimeIndex} from the start
timestamp and sample rate:
\[
t_i = t_{\text{start}} + \frac{i}{f_s},\qquad i = 0,1,\dots,N-1 .
\]
The ACC reader divides raw counts by 64 to obtain $g$. The IBI reader places
each beat at $t_{\text{start}}+\text{offset}_s$, producing an irregularly-sampled
series. All reads are wrapped by \code{\_safe\_read}: a missing or malformed file
yields an empty Series/DataFrame plus a warning, so the pipeline degrades
gracefully to partial data instead of crashing. The loader also auto-resolves
both flat (\code{SXX/HR.csv}) and nested (\code{<SID>\_E4\_Data/HR.csv}) layouts.

%=======================================================================
\section{Synchronization and Resampling}\label{sec:sync}
%=======================================================================

\code{src/data/sync.py::build\_synced\_frame} aligns all modalities onto a common
1\,Hz grid (\code{cfg.resample.freq = "1s"}). 1\,Hz is chosen because HR is
natively $\sim$1\,Hz, it is fine enough for minute-scale forecasting, and it is
cheap enough for on-device streaming. Per-signal handling:

\begin{itemize}[leftmargin=1.4em,itemsep=2pt]
\item \textbf{HR}: clipped to a plausible range, then \code{resample("1s").mean()}
      to absorb tiny timestamp drift.
\item \textbf{ACC (32\,Hz)}: gravity removed (Section~\ref{sec:clean}), then the
      per-axis magnitude
      $\lVert a \rVert = \sqrt{a_x^2 + a_y^2 + a_z^2}$
      is aggregated per second into \textbf{mean}, \textbf{std}, and \textbf{max}
      (\code{ACC\_mag\_mean/std/max}) --- capturing both average intensity and
      burstiness within the second.
\item \textbf{EDA (4\,Hz)}: low-pass smoothed, then per-second mean.
\item \textbf{TEMP (4\,Hz)}: per-second mean (slow, clean signal).
\item \textbf{BVP (64\,Hz)}: band-pass cleaned but \emph{kept at native rate}, so
      peak detection has full temporal resolution.
\item \textbf{IBI (irregular)}: implausible beats dropped, kept on its native
      irregular index for beat-level HRV.
\end{itemize}

The four grid-aligned columns (HR, EDA, TEMP, ACC aggregates) are concatenated,
sorted, and passed through short-gap interpolation. The function returns a dict
with \code{grid}, native-rate \code{bvp\_clean}, irregular \code{ibi\_clean}, and
the \code{tags} index.

%=======================================================================
\section{Signal Cleaning and Filtering}\label{sec:clean}
%=======================================================================

Cleaning is signal-specific (\code{src/data/cleaning.py}). All frequency-domain
filters are zero-phase Butterworth filters applied with
\code{scipy.signal.filtfilt} (forward-backward), which doubles the effective
order and, crucially, introduces \emph{no phase delay} --- essential when the
timing of pulse peaks and HR changes carries physiological meaning.

\subsection{Band-pass (BVP / PPG)}
BVP is band-pass filtered to $[0.5, 8]\,$Hz (order 3) to remove baseline drift
below the pulse band and high-frequency noise above it. With Nyquist
$f_{\text{nyq}} = f_s/2$, the normalized cutoffs are
\[
W_{\text{low}} = \max\!\left(\frac{0.5}{f_{\text{nyq}}},\,10^{-4}\right),\qquad
W_{\text{high}} = \min\!\left(\frac{8}{f_{\text{nyq}}},\,0.999\right),
\]
and the biquad coefficients come from \code{butter(3, [W_low, W_high], "band")}.

\subsection{Low-pass (EDA / TEMP)}
EDA is low-pass filtered at 1\,Hz (order 3) before being split into tonic and
phasic components (Section~\ref{sec:features}). The normalized cutoff is
$W_c = \min(f_c/f_{\text{nyq}}, 0.999)$.

\subsection{Gravity removal (ACC)}
Each accelerometer axis is high-pass filtered at 0.5\,Hz to strip the DC gravity
component, leaving only \emph{dynamic} acceleration (movement) before the
magnitude is computed. This prevents posture (a static tilt) from registering as
activity.

\subsection{Range clipping and plausibility}
\begin{itemize}[leftmargin=1.4em,itemsep=2pt]
\item HR is clipped to $[30, 220]$\,bpm.
\item IBI is filtered to $[0.3, 2.0]\,$s (i.e.\ $30$--$200$\,bpm); values outside
      this are physiologically impossible and are dropped
      (\code{clean\_ibi}).
\end{itemize}

\subsection{Gap handling and segmentation}
\code{fill\_small\_gaps} linearly interpolates gaps up to
\code{max\_interpolate\_gap\_s} (5\,s). Each interpolated sample is tagged with a
\code{<col>\_stale} flag ($=1$ when the value was filled across a $>1$-sample
gap), so the model and the signal-quality layer can down-weight it. Gaps longer
than \code{max\_segment\_gap\_s} (60\,s) are treated as true discontinuities:
\code{segment\_by\_gaps} splits the recording into contiguous segments, and
sequence windows are never built across a segment boundary. This means a
$\sim$1-hour session broken by charging/removal yields several clean segments
rather than one corrupted one.

%=======================================================================
\section{Feature Engineering}\label{sec:features}
%=======================================================================

\code{src/features/build\_features.py::build\_feature\_table} orchestrates all
feature modules into one per-second table indexed by the 1\,Hz grid. Modules are
included conditionally on sensor availability, so a dataset lacking (say) EDA
simply omits EDA features. The assembly order matters because later modules
consume earlier ones:

\begin{enumerate}[leftmargin=1.4em,itemsep=1pt]
\item Base grid $+$ HR features $+$ IBI/HRV features $+$ time features (always).
\item ACC, BVP, TEMP, EDA features (if present).
\item Circadian baseline (needs HR).
\item Physiological state classifier (needs activity $+$ HR $+$ circadian).
\item Recovery features (need state).
\item Subject-relative baselines and long-context summaries.
\item Activity-gating features.
\item Forecast targets (future HR/RMSSD by shifting).
\end{enumerate}

\subsection{HR features (\texttt{hr\_features.py})}
For each window $w \in \{30, 60, 300\}$\,s over the HR series $h(t)$:
\[
\mu^{\text{HR}}_w(t)=\frac{1}{w}\!\!\sum_{\tau=t-w+1}^{t}\!\! h(\tau),\quad
\sigma^{\text{HR}}_w(t)=\sqrt{\tfrac{1}{w-1}\!\!\sum_{\tau=t-w+1}^{t}\!\!\big(h(\tau)-\mu^{\text{HR}}_w(t)\big)^2}.
\]
A trailing-window \textbf{trend} is the least-squares slope of $h$ over $w$
seconds; with $x=0,\dots,w-1$ and $\bar x$ its mean,
\[
\text{slope}_w(t)=\frac{\sum_x (x-\bar x)\big(h-\bar h\big)}{\sum_x (x-\bar x)^2}
\quad[\text{bpm/s}].
\]
Rate-of-change features $\text{HR\_roc}_\Delta(t)=h(t)-h(t-\Delta)$ are computed
for $\Delta \in \{5, 30, 60\}$\,s.

\subsection{IBI / HRV features (\texttt{ibi\_features.py})}
Wrist-PPG HRV is noisy, so beats are cleaned \emph{before} HRV is computed. A
\textbf{Hampel filter} rejects ectopic/artifact beats: with local median $m_i$
over a 5-beat centered window and median absolute deviation $\text{MAD}_i$, a
beat is flagged when
\[
\lvert b_i - m_i\rvert > n_\sigma \cdot 1.4826\cdot \text{MAD}_i,\qquad n_\sigma=3,
\]
where $1.4826$ makes the MAD a consistent estimator of the standard deviation
for Gaussian data. Flagged beats are replaced by the local median, and any
successive difference that spans a rejected beat is invalidated. Over each HRV
window $w\in\{60,300\}$\,s (requiring $\geq$\code{hrv\_min\_beats} valid beats):

\begin{align*}
\text{RMSSD}_w &= \sqrt{\tfrac{1}{M-1}\textstyle\sum_i \big(b_{i+1}-b_i\big)^2}
   &&\text{(short-term parasympathetic HRV)}\\
\text{SDNN}_w &= \operatorname{std}\big(\{b_i\}\big)
   &&\text{(overall HRV)}\\
\text{pNN50}_w &= \frac{1}{M}\sum_i \mathbb{1}\big[\lvert b_{i+1}-b_i\rvert > 50\text{ms}\big]
   &&\text{(fraction of large beat-to-beat jumps)}\\
\text{HR}^{\text{IBI}}_w &= \frac{60}{\overline{b}}
   &&\text{(HR implied by mean IBI)}
\end{align*}
Each beat-indexed metric is forward-filled onto the 1\,Hz grid. A quality column
\code{RMSSD\_valid\_fraction\_{w}s} (the rolling fraction of non-rejected beats)
is emitted for the signal-quality layer.

\subsection{Accelerometer features (\texttt{acc\_features.py})}
From \code{ACC\_mag\_mean} (gravity-removed magnitude in $g$), rolling intensity
and its std are computed over $w\in\{10,60\}$\,s. The primary intensity is
bucketed into an ordinal \textbf{activity level} using the thresholds
$\{0.05, 0.2, 0.5\}\,g$:
\[
\text{bucket} \in \{\text{REST}, \text{LIGHT}, \text{MODERATE}, \text{VIGOROUS}\},
\]
plus a binary \code{movement\_flag} $=\mathbb{1}[\text{intensity}>0.05\,g]$.

\subsection{BVP / pulse features (\texttt{bvp\_features.py})}
On the native-rate cleaned BVP, systolic peaks are found with
\code{scipy.signal.find\_peaks}, enforcing a minimum inter-peak distance of
$\lfloor f_s\cdot 60/220\rfloor$ samples (a 220\,bpm ceiling). Peak
\emph{prominences} give a pulse-amplitude series. Per second, the pulse count is
resampled and a rolling 60\,s sum yields \code{BVP\_pulse\_rate\_1min}; rolling
30\,s mean/std of amplitude give \code{BVP\_amplitude\_mean/std\_30s}.

\subsection{EDA features (\texttt{eda\_features.py})}
The tonic (slow) level is a 60\,s rolling mean; the phasic (fast, SCR) component
is the residual:
\[
\text{EDA}_{\text{tonic}}(t)=\tfrac{1}{60}\!\!\sum_{\tau=t-59}^{t}\!\!\text{EDA}(\tau),
\quad
\text{EDA}_{\text{phasic}}(t)=\text{EDA}(t)-\text{EDA}_{\text{tonic}}(t),
\]
with a rolling mean of $\lvert\text{phasic}\rvert$ as a 60\,s arousal proxy.

\subsection{Temperature features (\texttt{temp\_features.py})}
A 300\,s rolling mean and slope capture drift; deviation from a personal baseline
uses an \emph{expanding} mean (a proxy for the calibration-period baseline):
\[
\text{TEMP\_baseline\_deviation}(t)=\text{TEMP}(t)-\frac{1}{t}\sum_{\tau\le t}\text{TEMP}(\tau).
\]

\subsection{Time and circadian features (\texttt{time\_features.py}, \texttt{circadian.py})}
Time-of-day and day-of-week are encoded cyclically so midnight is adjacent to
23:59:
\[
\Big(\sin\tfrac{2\pi H}{24},\cos\tfrac{2\pi H}{24}\Big),\quad
\Big(\sin\tfrac{2\pi D}{7},\cos\tfrac{2\pi D}{7}\Big),
\]
where $H$ is fractional hour and $D$ day-of-week. The \textbf{causal circadian
baseline} answers ``what is normal HR for this person at this hour?'' For each
hour bucket $h\in\{0,\dots,23\}$ it maintains an expanding mean/std of HR,
\emph{shifted by one step} so the current sample never sees itself:
\[
z^{\text{circ}}(t)=\frac{\text{HR}(t)-\mu^{\text{circ}}_{h(t)}(t^-)}
{\max\!\big(\sigma^{\text{circ}}_{h(t)}(t^-),\,10^{-6}\big)}.
\]
On a short single session these are degenerate (one hour bucket), but the same
code scales directly to weeks of continuous data where each bucket accumulates
many prior days.

\subsection{Physiological state classifier (\texttt{state\_classifier.py})}
Each 1\,Hz sample is labeled with one of
\emph{sleep, rest, focused\_work, walking, exercise, recovery, unknown}. This
gives the anomaly layer context: an HR spike during \emph{exercise} is expected;
the same spike during \emph{sleep} is not. The decision cascade (all causal):
\begin{enumerate}[leftmargin=1.4em,itemsep=1pt]
\item activity $\geq$ MODERATE $\Rightarrow$ \textbf{exercise}.
\item recent exercise (600\,s lookback) and HR still $>$ circadian mean $+3$
      $\Rightarrow$ \textbf{recovery}.
\item activity $=$ LIGHT $\Rightarrow$ \textbf{walking}.
\item activity $=$ REST, night hours (23:00--05:59), HR $\le$ circadian mean
      $\Rightarrow$ \textbf{sleep}.
\item activity $=$ REST, daytime, elevated EDA arousal $\Rightarrow$
      \textbf{focused\_work}.
\item activity $=$ REST otherwise $\Rightarrow$ \textbf{rest}.
\end{enumerate}
Missing inputs degrade gracefully to \emph{unknown}.

\subsection{Recovery / fitness features (\texttt{recovery\_features.py})}
During and after an exercise bout the module tracks the bout's peak HR and how
fast HR falls from it. Heart-rate recovery is a well-known fitness marker:
\[
\text{HRR}(t)=\max(0,\ \text{HR}_{\text{peak}}-\text{HR}(t)),\qquad
\text{HRR rate}(t)=\frac{\text{HRR}(t)}{\Delta t_{\text{since peak}}/60}\ [\text{bpm/min}].
\]
A faster drop indicates better cardiovascular fitness; persistently slow
recovery over weeks flags declining fitness or incomplete recovery. All outputs
default to 0 outside a bout, so no NaNs are introduced.

\subsection{Subject-relative baselines and long context (\texttt{subject\_baseline.py})}
Absolute HR/TEMP/RMSSD vary enormously between people, so deviation from a
person's \emph{own} recent baseline is more informative than the raw value. For
each signal $s$ and baseline window $W=3600$\,s (causal):
\[
s_{\text{minus\_baseline}}(t)=s(t)-\mu^s_W(t),\qquad
s_{\text{z\_rolling}}(t)=\frac{s(t)-\mu^s_W(t)}{\sigma^s_W(t)}.
\]
Multi-scale long-context summaries of HR and RMSSD over $W\in\{900,1800,3600\}$\,s
(15/30/60\,min) provide \code{ctx\_mean}, \code{ctx\_std}, and an endpoint slope
$\big(s(t)-s(t-W)\big)/W$, giving the model both the immediate trajectory and the
slow baseline trend without needing a 60-minute raw input window. All deviation
columns default to 0 (never NaN) so windowing never drops rows.

\subsection{Activity-gating features (\texttt{activity\_gating.py})}
``High HR'' means different things at different motion levels: HR-high $+$
ACC-high is likely exercise; HR-high $+$ ACC-low is more suspicious (stress,
illness, arousal). For each activity bucket, a strictly-causal running mean/variance
of HR is maintained with \textbf{Welford's algorithm}:
\[
\mu_k=\mu_{k-1}+\frac{x_k-\mu_{k-1}}{k},\qquad
M_{2,k}=M_{2,k-1}+(x_k-\mu_{k-1})(x_k-\mu_k),\qquad
\sigma_k^2=\frac{M_{2,k}}{k}.
\]
This yields \code{hr\_expected\_for\_activity}, the residual, a clipped
\code{hr\_z\_for\_activity} $\in[-8,8]$, and an explicit
\code{hr\_high\_while\_still} flag (low motion but elevated HR-for-activity).
NaN HR (sensor gaps) never enters the running stats, so a gap cannot poison the
baseline.

\subsection{Feature selection and forecast targets}
\code{numeric\_feature\_columns} selects the model inputs: it keeps every numeric
column and drops (i) target columns (any \code{*\_target\_*}) and (ii) the
categorical name columns (\code{activity\_bucket\_name},
\code{physio\_state\_name}). Forecast \textbf{targets} are built by shifting HR
(and, when \code{predict\_rmssd} is on, RMSSD) backward by each horizon:
\[
\text{HR\_target}_h(t)=\text{HR}(t+h),\qquad h\in\{60,300,600\}\text{ s}.
\]
Optional targets (config flags) add change-from-now deltas
($\text{HR}(t+h)-\text{HR}(t)$) and other vitals (TEMP/EDA) for multi-task
learning.

%=======================================================================
\section{Windowing and Dataset Construction}\label{sec:windows}
%=======================================================================

\code{src/models/datasets.py::make\_windows} slides a length-$L$ window
(\code{input\_seq\_len\_s} $=120$\,s) over the feature table with a stride
(\code{stride\_s} $=5$\,s), producing
\[
X\in\mathbb{R}^{N\times L\times F},\quad
y\in\mathbb{R}^{N\times H},\quad
\text{end\_idx}\in\mathbb{Z}^{N},
\]
where $F$ is the number of features and $H$ the number of horizons. A window is
\emph{dropped} if any feature or target within it is NaN --- this is exactly why
so many features are engineered to default to 0 rather than NaN. The short
120\,s context is deliberate: it keeps fragmented $\sim$1-hour sessions
productive while the long-context features (above) carry the slow information a
longer raw window would otherwise be needed for.

\textbf{Chronological split.} Data is split forward in time (never shuffled
across time) into train/val/test by \code{val\_fraction} and
\code{test\_fraction} ($0.15$ each), so evaluation always predicts the future
from the past.

\textbf{Scaling.} A \code{StandardScaler} is fit on the flattened training
windows $(N\!\cdot\!L, F)$ and applied to all splits. Targets are
\emph{also} standardized (\code{fit\_target\_scaler}), because HR (bpm) and RMSSD
(ms) live on very different scales and mixing them raw destabilizes the
Gaussian-NLL objective. Predictions are inverse-transformed back to physical
units for reporting; the predicted std is rescaled by the per-target scale,
$\sigma_{\text{raw}}=\sigma_{\text{scaled}}\cdot\text{scale}$.

%=======================================================================
\section{Forecasting Models}\label{sec:models}
%=======================================================================

All sequence models share a \textbf{heteroscedastic head}: from the
representation at the most recent timestep they emit both a mean and a positive
standard deviation per horizon,
\[
\hat\mu_h,\ \hat\sigma_h=\text{softplus}(\cdot)+10^{-3},
\]
so every forecast comes with its own uncertainty --- the quantity the anomaly
layer normalizes by.

\subsection{Temporal Convolutional Network (primary)}
The recommended model is a causal dilated TCN (\code{src/models/tcn.py}). Each
\code{CausalConvBlock} is a 1-D convolution with left-padding
$(k-1)\cdot d$ that is trimmed on the right (\emph{causal trim}), a ReLU,
dropout, and a residual connection:
\[
\text{out}=\text{Dropout}\big(\text{ReLU}(\text{CausalConv}_{k,d}(x))\big)+\text{Res}(x).
\]
Stacking $\ell=0,\dots,L{-}1$ levels with exponential dilation $d=2^\ell$ gives an
exponentially growing receptive field
\[
\text{RF}=1+\sum_{\ell=0}^{L-1}2(k-1)2^{\ell}.
\]
With the tuned config ($k=3$, levels $=4$) the RF is $\approx$61\,s, matched to the
120\,s window --- an earlier $k=5$, 6-level version had RF~$\approx$505\,s
$\gg$ window and overfit. Hidden width is 32 (down from 64 to avoid overfitting
on a few hundred windows/subject).

\subsection{Recurrent baselines (LSTM / GRU)}
\code{src/models/lstm\_gru.py} provides an LSTM and a GRU sharing the same
mean$+$std head, reading the sequence and predicting from the last hidden state.

\subsection{Transformer}
\code{src/models/transformer.py} projects features to $d_{\text{model}}$, adds
sinusoidal positional encoding, and applies a causal-masked
\code{TransformerEncoder} so predictions depend only on past/current steps.

\subsection{XGBoost (tabular baseline)}
\code{src/models/xgb\_model.py} trains one gradient-boosted regressor per target
on a single feature vector per timestep (the engineered rolling features already
encode recent history), so it needs no sequence window. Uncertainty is
approximated by the constant per-target residual std on the validation set.

\subsection{$\beta$-NLL loss and training}
Models are trained (\code{src/models/train.py}) by minimizing the
\textbf{$\beta$-NLL} objective (Seitzer et al., 2022):
\[
\mathcal{L}=\frac{1}{NH}\sum_{i,h}
\Big(\underbrace{\tfrac12\log(2\pi\hat\sigma_{ih}^2)+\tfrac{(y_{ih}-\hat\mu_{ih})^2}{2\hat\sigma_{ih}^2}}_{\text{Gaussian NLL}}\Big)\cdot
\big(\hat\sigma_{ih}^2\big)^{\beta}_{\text{stop-grad}},\qquad \beta=0.5.
\]
Plain Gaussian NLL can be minimized by \emph{inflating} the variance (predicting
the mean with huge $\sigma$), because the $1/\sigma^2$ term suppresses the error
--- the ``variance collapse'' that makes forecasts flat. Multiplying by
$\text{stop-grad}(\sigma^2)^\beta$ restores the mean's gradient magnitude;
$\beta{=}0$ recovers standard NLL, $\beta{=}1$ is MSE on the mean. Optimization
uses Adam; \textbf{model selection is by validation MSE of the mean}, not NLL, so
a variance-collapsed model can never appear ``best.'' Forecast quality is
reported per horizon as MAE, RMSE, and MAPE.

%=======================================================================
\section{Personalization}\label{sec:personal}
%=======================================================================

The system learns \emph{per person} through a staged transfer path
(\code{src/pipeline/}).

\begin{longtable}{@{}p{0.26\textwidth}p{0.68\textwidth}@{}}
\toprule
\textbf{Stage} & \textbf{What it does} \\
\midrule
\endhead
Global pretraining (\code{pretrain\_pipeline.py}) & Pools windows from many
subjects, intersects their feature columns, trains one population model. Serves
every new user until they have enough personal data. \\
SSL warm-start (\code{ssl\_pretrain\_pipeline.py}) & Self-supervised masked
reconstruction: randomly mask $25\%$ of feature values and train a TCN backbone
to reconstruct them (no labels). The learned backbone warm-starts a forecaster,
helping subjects with little data. \\
Per-user fine-tune (\code{finetune\_pipeline.py}) & Continues training the
pretrained model on one subject's windows (optionally freezing the backbone),
in the pretrained target-scaler space. Saved as a normal loadable checkpoint. \\
Online update (\code{online\_update\_pipeline.py}) & Periodically (e.g.\ nightly)
continues training on the \emph{most recent} windows only, at a small LR for a
few epochs, and \code{partial\_fit}s the feature scaler --- so the model tracks
slow shifts (e.g.\ improving fitness lowering resting HR over months). \\
Multi-dataset harness (\code{multi\_dataset\_pipeline.py}) & Leave-one-dataset-out
and repeated-split evaluation across pooled datasets. \\
\bottomrule
\end{longtable}

\subsection{Personal Digital Twin (\texttt{personalization/digital\_twin.py})}
The twin is a compact JSON profile of a person's own physiology: resting, sleep,
walking, and running HR (means per physiological state), average RMSSD, and a
per-hour-of-day circadian HR table. The \textbf{health-drift score} compares the
live HR to this personal expectation:
\[
\text{DigitalTwinScore}(t)=\frac{\text{HR}(t)-\text{ExpectedHR}(t)}{\text{ExpectedStd}(t)},
\]
where $\text{ExpectedHR}(t)$ is read from the circadian table for the current
hour, falling back to the state-specific baseline (sleep/rest/walk/run) when no
circadian estimate exists yet. This personal baseline --- not a population norm
--- is the reference for drift.

\subsection{Kalman baseline tracker (\texttt{models/kalman.py})}
A 1-D constant-velocity Kalman filter tracks a slowly-varying scalar (resting HR,
or the forecast residual) with state $[\text{level},\text{trend}]$:
\[
x_k=Fx_{k-1}+w,\ w\sim\mathcal N(0,Q);\qquad z_k=Hx_k+v,\ v\sim\mathcal N(0,R).
\]
The normalized innovation $\tilde y_k=(z_k-H\hat x_{k}^-)/\sqrt{S_k}$ is used by
the anomaly layer to catch gradual drift.

%=======================================================================
\section{Anomaly Detection}\label{sec:anomaly}
%=======================================================================

The anomaly layer (\code{src/anomaly/}) is deliberately \emph{separate} from the
forecaster: it consumes the model's per-timestep outputs
$(y_t, \hat\mu_t, \hat\sigma_t)$ and turns them into an interpretable score and
alerts, saving every intermediate quantity for debugging.

\subsection{Uncertainty-normalized residuals}
The core signal is the residual normalized by the model's own predicted
uncertainty (\code{detector.py}):
\[
z_h(t)=\frac{\lvert y_h(t)-\hat\mu_h(t)\rvert}{\hat\sigma_h(t)+\varepsilon},
\quad \text{clipped to }[0,8].
\]
Clipping prevents a single artifact from exploding the score. Per-target
$z$-scores are blended with separate HR/RMSSD weights and a mild shorter-horizon
emphasis.

\subsection{Adaptive (EWMA) baselines with warm-up (\texttt{adaptive.py})}
``Normal'' is contextual and drifts, so residual statistics are tracked with an
exponentially-weighted mean/variance, optionally \emph{per activity bucket}:
\[
\mu \leftarrow \mu + \lambda(x-\mu),\qquad
\sigma^2 \leftarrow (1-\lambda)\big(\sigma^2+\lambda(x-\mu)^2\big),\qquad \lambda=0.02.
\]
Variance is initialized large ($\sigma\!\approx\!10$\,bpm) and $z$-scores are
suppressed to 0 for the first 50 observations per key, killing cold-start false
alarms.

\subsection{Multi-channel evidence}
The scorer (\code{scoring.py}) combines complementary detectors, each catching a
different failure mode:
\begin{itemize}[leftmargin=1.4em,itemsep=2pt]
\item \textbf{Forecast $z$-scores} --- point deviations from the personal
      forecast, per horizon.
\item \textbf{Prediction-interval score} $\lvert y-\hat\mu\rvert/\hat\sigma$ and
      $\pm1.96\hat\sigma$ out-of-interval flags.
\item \textbf{RMSSD drop} --- a one-sided $z$-score and fractional drop of HRV
      below its EWMA baseline (HRV collapse).
\item \textbf{Kalman innovation} --- normalized innovation of the 1-min residual
      stream (slow drift).
\item \textbf{CUSUM drift} --- a two-sided decaying CUSUM on residuals that
      accumulates sustained directional bias per-step $z$-scores normalize away:
      \[
      s^+_i=\max\!\big(0,\ \text{decay}\cdot s^+_{i-1}+z_i-k\big),\quad
      s^-_i=\max\!\big(0,\ \text{decay}\cdot s^-_{i-1}-z_i-k\big),
      \]
      with slack $k=0.7$ and $\text{decay}=0.998$ ($\approx$6-min half-life), so
      exercise bursts do not bleed into later rest.
\item \textbf{Circadian $z$} and \textbf{digital-twin $z$} --- deviation from the
      person's time-of-day and state baselines.
\end{itemize}

\subsection{Illness composite (\texttt{explain.py})}
Three signals classically co-move during infection/fever; individually each may
be noise, together they are strong evidence:
\[
\text{illness}=\frac13\!\left(
\frac{[z^{\text{HR}}]_+}{3}
+\frac{[\text{RMSSD}_{\%\text{drop}}]_+}{0.35}
+\frac{[\Delta\text{TEMP}]_+}{1.0}\right),
\]
scaled so a ``typical fever'' (HR $z\!\approx\!3$, RMSSD drop $\approx$35\%, TEMP
$+1^\circ$C) scores $\approx 1.0$.

\subsection{Combined score, exercise suppression, illness boost}
\[
\text{score}(t)=\max_{\text{channels}}\Big(w_{\text{sup}}(t)\cdot\text{channel}\Big)
+1.5\cdot[\text{illness}(t)]_+ .
\]
During \emph{exercise/recovery} the HR-based channels (forecast, circadian,
digital-twin, Kalman, CUSUM) are down-weighted by $80\%$
($w_{\text{sup}}=0.2$) because elevated HR is expected; RMSSD and illness terms
are \emph{not} suppressed. The base score is a $\max$ over channels so one strong
signal fires immediately, while the additive illness term lets several
individually-mild channels jointly cross the threshold.

\subsection{Signal-quality gating (\texttt{signal\_quality.py})}
Real wrist data is full of motion/contact artifacts; without a quality gate,
sensor noise is mistaken for a health event. A per-second Signal-Quality Index
$\in[0,1]$ is computed per modality (implausible jumps, out-of-range values,
stale/interpolated flags, HRV valid-fraction, ACC saturation, BVP flatline). The
overall SQI is the \emph{weakest} modality (a single bad sensor spoils the
window), and the effective anomaly score is
\[
\text{score}_{\text{eff}}(t)=\text{score}(t)\cdot \text{SQI}_{\text{overall}}(t),
\]
so unreliable windows are labeled ``sensor issue,'' not ``anomaly.''

\subsection{Hysteresis alerting and threshold calibration (\texttt{detector.py})}
Alerts use a start/stop hysteresis state machine to avoid chatter:
\begin{itemize}[leftmargin=1.4em,itemsep=2pt]
\item An alert \textbf{starts} only after $\text{score}_{\text{eff}}\ge
      \theta_{\text{start}}$ for \code{min\_duration\_s} continuous seconds
      (the build-up is backfilled).
\item It \textbf{stays} until the score drops below
      $\theta_{\text{stop}}=0.6\,\theta_{\text{start}}$.
\item After it ends, a \code{cooldown\_s} quiet period applies, bypassed only for
      genuine emergencies ($\ge 1.5\,\theta_{\text{start}}$).
\end{itemize}
The start threshold is \textbf{calibrated} on the distribution of
\emph{normal} (non-anomalous) validation scores at a chosen percentile
($99$th, floored at 3.0), targeting a false-alert budget rather than a
hand-picked number. \code{alert\_stats} reports the normal alert rate,
alerts-per-hour, and mean alert duration; \code{severity} maps the continuous
score to \emph{normal / watch / alert / urgent}. In the served AI path,
severity is bucketed as \emph{normal} $<1.5\le$ \emph{watch} $<3.0\le$
\emph{alert}.

\subsection{Explainable reasons}
Every alert carries a human-readable reason built from the same numbers, e.g.\
\emph{``+HR 22\,bpm above expected, RMSSD dropped 35\%, Temp elevated
$+1.2^\circ$C, Activity: rest''}, so a clinician sees \emph{why}, not just a
score.

%=======================================================================
\section{Simulation and Evaluation}
%=======================================================================

\code{src/simulation/} generates synthetic 1\,Hz traces (circadian HR baseline
$+$ activity-driven response) with injected abnormal windows labeled by an
\code{anomaly\_label} column, and can concatenate scenarios onto a real subject's
grid for ``inject-into-real-data'' testing. Evaluation scripts
(\code{scripts/run\_*.py}) cover preprocessing, training, inference, ablation,
error analysis, online update, and simulation scoring. Forecast metrics are
per-horizon MAE/RMSE/MAPE; calibration (\code{evaluation/calibration.py})
checks that predicted uncertainty matches observed error; the anomaly harness
reports detection rate and false-alarms-per-hour.

%=======================================================================
\section{AI Serving Layer}\label{sec:aiservice}
%=======================================================================

The \code{ai/} FastAPI service exposes the research pipeline over HTTP with a
stable contract and a never-fail fallback.

\subsection{Model resolution (\texttt{hrv\_forecast/model\_loader.py})}
\code{HRV\_MODEL\_DIR} is treated as a \emph{base} directory. Per request, a
checkpoint is resolved in priority order: the subject's personal directory
$\rightarrow$ \code{global/} $\rightarrow$ the base directory as a flat
checkpoint. Each resolved directory is loaded once and cached, so many subjects
sharing the global model pay the load cost only once. Any load error is logged
and the service falls back to the mock engine rather than crashing.

\subsection{Feature adapter (\texttt{hrv\_forecast/feature\_adapter.py})}
The adapter turns an API list of \code{HRVSample}s into the exact feature table
the research pipeline expects --- without raw E4 files --- by building the
``synced'' shape (grid, cleaned IBI) directly and calling the real
\code{build\_feature\_table}. The HR grid is built by resampling first and then
reindexing onto the 1\,Hz bins (a fix ensuring sub-second-timestamp inputs no
longer produce an all-NaN HR column).

\subsection{Real vs.\ mock engines and graceful fallback}
\code{real\_engine.py} runs the true pipeline: anomaly scoring and the digital
twin need only feature engineering (no trained net) and run for real today;
forecasting needs a trained checkpoint and raises if none is loaded.
\code{service.py} wraps each call so that if the real path raises for any reason
(too few samples, a bad checkpoint, an unexpected shape), it logs and returns the
deterministic mock engine's response --- callers always learn which they got via
\code{model\_status}. Endpoints (\code{routes/hrv.py}):
\code{GET /hrv/status}, \code{POST /hrv/forecast}, \code{POST /hrv/anomaly},
\code{POST /hrv/digital-twin}; plus \code{/health} and the agentic \code{/chat}.

\subsection{Agentic analysis workflow (\texttt{ai/workflow.py})}
A LangGraph state machine chains five nodes:
\[
\text{planner}\to\text{signal}\to\text{medical}\to\text{trend}\to\text{report}\to\text{END}.
\]
\emph{Planner} classifies intent (stress/sleep/hrv/trend/general) by keyword;
\emph{signal} runs the PPG analysis pipeline; \emph{medical} composes only
data-grounded reasoning (never invented numbers); \emph{trend} computes a real
first-half-vs-second-half comparison; \emph{report} calls the LLM
(\code{langchain-xai}, Grok) with a strict ``use only measured data'' system
prompt and falls back to a deterministic template if the LLM is unavailable.

\subsection{Standalone PPG analysis (\texttt{ai/analysis/})}
\code{HealthPipeline} preprocesses PPG (band-pass $0.5$--$8$\,Hz, normalize,
Savitzky--Golay smooth), detects peaks (\code{find\_peaks}, min distance
$0.45\,$s), forms RR intervals, and computes HRV time-domain metrics
(RMSSD, SDNN, mean RR, mean HR, pNN50) and a heuristic stress score
$\text{clip}\big([60-\text{RMSSD}]_+ + [\text{HR}-70]_+,\,0,100\big)$ mapped to
Low/Moderate/High. Numeric results are cast to native Python floats so NumPy
scalar reprs never leak into reports.

%=======================================================================
\section{Backend Platform (Django REST)}
%=======================================================================

The \code{backend/} app (\code{login/}) is API-only (DRF, JWT via SimpleJWT),
grouped under \code{/api/} in \code{mHealth/urls.py}.

\subsection{Authentication}
Email-OTP signup and login with JWT access/refresh tokens:
\code{/api/auth/signup}, \code{send-otp}, \code{verify-otp}, \code{login},
\code{login/verify}, \code{logout}, \code{change-password}, \code{me},
\code{token/refresh}. \code{DJANGO\_SECRET\_KEY} is environment-only and startup
fails loudly in production if unset; \code{DEBUG} defaults to \emph{False}
(fail-closed). Settings now load \code{backend/.env} explicitly via
\code{load\_dotenv} so hand-launched servers honor it.

\subsection{Data model (\texttt{login/models.py})}
\begin{longtable}{@{}p{0.28\textwidth}p{0.66\textwidth}@{}}
\toprule
\textbf{Model} & \textbf{Purpose} \\
\midrule
\endhead
\code{UserProfile} & Extended per-user profile. \\
\code{EmailOTP} & One-time passcodes for signup/login. \\
\code{Device} & Registered ingestion device (key, battery, firmware, session). \\
\code{GoogleSheet}, \code{ExcelFile} & External/tabular data sources. \\
\code{UploadedDocument} & Category metadata (lab/prescription/imaging/\ldots) for
files on disk. \\
\code{HRVAnomalyAlert} & A fired anomaly/digital-twin deviation from
\code{/hrv/anomaly} (score, reasons, severity, emailed flag). \\
\code{AlertRule}, \code{AlertFired} & Per-user threshold rules and their firings. \\
\code{SessionAnnotation} & Timestamped notes on a session. \\
\code{ContactMessage}, \code{ParticipantConsent} & Contact form, health consent. \\
\bottomrule
\end{longtable}

\subsection{Signal Portal and HRV in the backend}
The portal (\code{/api/device/sessions/*}) aggregates E4 session folders:
per-minute/hour/day stats for HR/EDA/TEMP, ACC movement magnitude, a BVP
waveform sample, and HRV from IBI. Backend HRV mirrors the research formulas:
$\text{RMSSD}=\sqrt{\overline{(\Delta b)^2}}$ and $\text{SDNN}=\operatorname{std}(b)$
(\code{views/device.py::\_hrv\_stats}). \code{session\_trends} and
\code{compare\_sessions} power the trends/compare pages.

\subsection{HL7 export and alerts}
\code{/api/hl7/*} converts session CSVs into HL7 messages (text and PDF via
ReportLab). \code{/api/alerts/*} manages threshold rules, evaluates them when
a session loads, records firings, and ingests HRV anomalies from the AI service
(\code{report\_hrv\_anomaly}), emailing the user on \emph{alert} severity.

%=======================================================================
\section{Frontend and Device Ingestion}\label{sec:frontend}
%=======================================================================

\subsection{Next.js web app}
The App-Router frontend (\code{frontend/src/app/}) provides pages for
\emph{login, dashboard, files, PPG visualization, portal
(+session/trends/compare), alerts, devices, consent}, and \emph{admin}
(overview, patients). \code{src/lib/api.ts} is an Axios client with a 15\,s
timeout that attaches the JWT and transparently refreshes on 401.
\code{HRVInsights.tsx} calls the AI service directly (with an 8\,s timeout) to
render forecasts, the digital twin, and anomaly badges alongside the recharts
time series. Styling uses the emerald ``CareVibe'' theme.

\subsection{Device heartbeat}
\code{device-ingestion/heartbeat.py} posts \code{battery/firmware/session} to
\code{/api/devices/heartbeat/} with an \code{X-Device-Key} header, so the portal
can show a device as online with live status. The key is issued once in the
portal's ``Register device'' flow.

%=======================================================================
\section{Configuration Reference}
%=======================================================================

Key parameters (\code{configs/config.yaml}); the AI feature adapter mirrors the
relevant subset as a literal dict so \code{ai/} needs no YAML dependency.

\begin{longtable}{@{}p{0.32\textwidth}p{0.20\textwidth}p{0.40\textwidth}@{}}
\toprule
\textbf{Parameter} & \textbf{Value} & \textbf{Meaning} \\
\midrule
\endhead
\code{resample.freq} & \code{1s} & common grid frequency \\
\code{hr\_min/max\_bpm} & 30 / 220 & HR clip range \\
\code{ibi\_min/max\_s} & 0.3 / 2.0 & plausible IBI range \\
\code{bvp\_bandpass\_hz} & [0.5, 8.0] & BVP band-pass \\
\code{eda\_lowpass\_hz} & 1.0 & EDA low-pass \\
\code{max\_interpolate\_gap\_s} & 5 & max gap to interpolate \\
\code{max\_segment\_gap\_s} & 60 & gap that splits a recording \\
\code{hr\_windows\_s} & 30/60/300 & HR rolling windows \\
\code{hrv\_windows\_s} & 60/300 & HRV rolling windows \\
\code{hrv\_min\_beats} & 3 & min valid beats per HRV window \\
\code{hrv\_hampel\_window/sigma} & 5 / 3.0 & ectopic-beat rejection \\
\code{baseline\_windows\_s} & 3600 & personal baseline window \\
\code{longctx\_windows\_s} & 900/1800/3600 & long-context summaries \\
\code{input\_seq\_len\_s} & 120 & model input window \\
\code{horizons\_s} & 60/300/600 & forecast horizons \\
\code{stride\_s} & 5 & window stride \\
\code{hidden\_channels} & 32 & model width \\
\code{tcn\_levels / kernel\_size} & 4 / 3 & TCN depth/kernel (RF $\approx$61\,s) \\
\code{dropout} & 0.2 & regularization \\
\code{epochs / lr / batch} & 40 / 1e-3 / 64 & training \\
\code{nll\_beta} & 0.5 & $\beta$-NLL weight \\
\code{val/test\_fraction} & 0.15 / 0.15 & chronological split \\
\code{ewma\_lambda} & 0.02 & residual EWMA rate \\
\code{z\_thresh} & 3.0 & base alert threshold \\
\code{start\_percentile} & 99.0 & calibration percentile \\
\code{stop\_fraction} & 0.6 & hysteresis stop ratio \\
\code{min\_duration\_s / cooldown\_s} & 20 / 60 & alert persistence/quiet \\
\bottomrule
\end{longtable}

%=======================================================================
\section{Reliability and Hardening Notes}
%=======================================================================

Several robustness properties are worth calling out because they shape the
working code throughout:
\begin{itemize}[leftmargin=1.4em,itemsep=2pt]
\item \textbf{Never-NaN features.} Deviation/recovery/context features default to
      0 rather than NaN so \code{make\_windows} does not silently discard rows.
\item \textbf{Strict causality.} Circadian, activity-gating, and EWMA baselines
      use expanding means with a one-step shift or online updates, so no feature
      sees the future.
\item \textbf{Fail-open service.} The AI HRV endpoints always answer: real
      pipeline first, deterministic mock on any exception, with the mode exposed
      via \code{model\_status}.
\item \textbf{Fail-closed security.} \code{DEBUG} defaults to False and the
      secret key must be supplied in production or startup aborts.
\item \textbf{Safe checkpoint loading.} \code{torch.load(..., weights_only=True)} is used when loading model/SSL state to avoid
      arbitrary-code-execution from pickled checkpoints.
\end{itemize}

\vfill
\noindent\rule{\linewidth}{0.4pt}\\[4pt]
\small This report documents the Care-Sync/M-Health codebase as implemented:
the data path (ingest $\to$ sync $\to$ clean $\to$ features $\to$ windows),
the forecasting models and their $\beta$-NLL training, the personalization
stack and digital twin, the multi-channel explainable anomaly detector, and the
FastAPI/Django/Next.js platform that serves it.
\end{document}
